\section{Data Management and Indexing}

\subsection{Raw Data Storage Layer}
\begin{itemize}
    \item \textbf{Raster Data:} All satellite imagery (HLS scenes) is stored as Cloud Optimized GeoTIFFs (COGs).
    \begin{itemize}
        \item Located in cloud object storage (AWS S3).
        \item COGs are optimized for fast, partial reads and cloud-native processing.
    \end{itemize}
    \item \textbf{Metadata:} Stored separately in a PostgreSQL database with PostGIS for efficient querying.
    \begin{itemize}
        \item Includes spatial (bounding box), temporal (acquisition date), and descriptive (band, cloud cover) attributes.
    \end{itemize}
\end{itemize}

\subsection{Ingestion and Catalog Building}

\subsection{Batch Ingestion (ETL Pipeline)}
A PySpark + Rasterio pipeline processes raw HLS files. For each image:
\begin{enumerate}
    \item Traverse directories or cloud buckets.
    \item Extract metadata (date, tile ID, band, cloud cover, bounding box).
    \item Convert to COG format if needed.
    \item Store metadata in PostgreSQL.
\end{enumerate}

\subsection{Spatial Indexing (PostGIS)}
Bounding box is read from each COG using Rasterio.
\begin{itemize}
    \item A geometry is created using \texttt{ST\_MakeEnvelope(...)} in PostGIS.
    \item A spatial index is built using \texttt{CREATE INDEX ... USING GIST}.
    \item Enables fast spatial queries like \texttt{ST\_Intersects} and \texttt{ST\_Contains}.
\end{itemize}

\section{Spatiotemporal Data Cube Index}
Final output: a centralized metadata table that links:
\begin{itemize}
    \item \textbf{Space} $\rightarrow$ Bounding box geometry.
    \item \textbf{Time} $\rightarrow$ Acquisition date.
    \item \textbf{Link} $\rightarrow$ File path or S3 URL to the COG.
\end{itemize}
Stored in PostgreSQL with PostGIS indexing. Queried using SQL for spatiotemporal filtering. This centralized index serves as an entry point for the downstream data transformation pipeline, enabling the efficient retrieval of filtered pixel-level reflectance data for feature computation as described in Section \ref{sec:data_transformation}.
